\documentclass[12pt]{article}
\newcommand{\DocShort}{Plan maestro}
\input{preamble_population.tex}
\begin{document}
\doccover{Método del plan maestro}{2026-05}
\handoutintro{Estimar la población futura sumando la capacidad poblacional de las zonas definidas en el plan urbano.}{Ciudades con instrumentos formales de ordenamiento territorial y zonificación.}

\section{Datos ejemplo}
Se requiere:
\begin{itemize}
\item zonificación del plan maestro,
\item área urbanizable por zona,
\item densidad adoptada para cada uso,
\item porcentaje de ocupación esperado en el horizonte de diseño.
\end{itemize}

\section{Ecuaciones mínimas}
La población de cada zona se estima como:
\begin{equation}
P_i=A_iD_if_i
\end{equation}
donde $P_i$ es la población de la zona $i$, $A_i$ es el área útil de la zona, $D_i$ es la densidad adoptada y $f_i$ es la fracción de ocupación proyectada.

La población total es:
\begin{equation}
P_T=\sum_i P_i
\end{equation}
donde $P_T$ es la población total de diseño del plan maestro.

\section{Procedimiento}
\begin{enumerate}
\item Identificar las zonas del plan maestro.
\item Asignar el área útil de cada zona.
\item Definir la densidad de diseño.
\item Corregir por el nivel de ocupación esperado.
\item Sumar la población de todas las zonas.
\end{enumerate}

\section{Ejemplo resuelto}
Supóngase el siguiente esquema de expansión urbana para el horizonte 2045.

\begin{center}
\begin{tabular}{lrrrr}
\toprule
Zona & Área (ha) & Densidad (hab/ha) & Ocupación & Población \\
\midrule
Residencial consolidada & 320 & 180 & 1.00 & 57,600 \\
Residencial de expansión & 210 & 160 & 0.85 & 28,560 \\
Residencial mixta & 95 & 220 & 0.90 & 18,810 \\
Área universitaria & 35 & 90 & 0.75 & 2,363 \\
\makecell[l]{Parque industrial\\con vivienda} & 60 & 110 & 0.70 & 4,620 \\
Centro de servicios & 28 & 140 & 0.80 & 3,136 \\
\midrule
Total & & & & 115,089 \\
\bottomrule
\end{tabular}
\end{center}

Por ejemplo, para la zona residencial de expansión:
\[
P = 210 \times 160 \times 0.85 = \num{28560}
\]

La población total de diseño es:
\[
P_T = \num{57600}+\num{28560}+\num{18810}+\num{2363}+\num{4620}+\num{3136}=\num{115089}
\]

\resultbox{La población de diseño estimada mediante el plan maestro es 115,089 habitantes.}

\section{Teoría breve}
Este método no extrapola directamente la historia censal. Parte de la estructura física y normativa de la ciudad futura, por eso es especialmente útil cuando la expansión está planificada.

\section{Observaciones de uso}
\begin{itemize}
\item Suele ser el método preferible cuando existe un plan urbano confiable.
\item Debe revisarse si las densidades normativas son realistas.
\item Puede combinarse con métodos históricos como verificación externa.
\end{itemize}
\end{document}
